% !TeX root = ../../Book2.tex
\chapter{一维波动方程}
\label{b2:ch:09}
\BookChapterNumber{b2:ch:09}
\begin{chapterabstract}
波动方程把局部振动与有限速度的传递联系起来。本章从特征变量推导 d'Alembert 公式，说明行波如何由初值确定，以及为什么特征线上的数据不足以确定解；再把非齐次项解释为各时刻速度增量的叠加，用收缩区间上的能量说明局部唯一性。全空间估计先从共同紧空间支集的解推广到有限能量解，再把演化延拓到初始位移属于 \(H^1\)、初始速度属于 \(L^2\)、外力属于 \(L^1_tL^2_x\) 的情形。行波、驻波、共振与折角传播将分别展示公式和估计所描述的不同侧面。
\end{chapterabstract}
\begin{learninggoals}
\begin{itemize}
\item\label{b2:item:09-goal-characteristic}从两个一阶传播量恢复位移，说明初始形状和初始速度各自决定什么，并解释特征初值线遗漏了哪一族数据。
\item\label{b2:item:09-goal-forcing}由 Duhamel 积分构造外力响应，在时间求导时保留积分上限的贡献。
\item\label{b2:item:09-goal-cone}画出依赖三角形，区分有限传播、支集外界与波前之后的尾部。
\item\label{b2:item:09-goal-energy}推导全局和局部能量恒等式，说明通量方向、边界运动和外力功怎样决定等式中的正负号，并用空间截断处理非紧支集的有限能量解。
\item\label{b2:item:09-goal-limit}区分经典解与能量解的数据和方程意义，依次完成光滑逼近、极限识别与能量解类内的唯一性证明。
\end{itemize}
\end{learninggoals}

把一根弦拉成一个隆起后静止释放，隆起会怎样分成两边传播的波？若在某一处短时间施力，后来一点会读取多大范围内的作用？波动方程同时涉及初始形状和初始速度，公式要区分二者；总能量又只控制速度和形状的梯度，不能单独控制位移本身。本章先把传播公式算出来，再用能量处理折角和仅有平方可积性的速度。

\ProseParagraphBreak

本章用到\chapref{b2:ch:01}的链式法则、含参变量积分与弱解的概念；第一编的 \(L^2\) 结构和磨光逼近用于最后的能量极限。本章将直接推导特征变量，因此不必预先学习一般特征理论。

\ProseParagraphBreak

\secref{b2:sec:0901}先求解齐次问题，随后在\subsecref{b2:subsec:090103}比较行波分解与初值线的选择；\secref{b2:sec:0902}把外力加入并解释传播区域。\secref{b2:sec:0903}先区分局部能量与全空间有限能量估计的条件，再构造低正则数据的能量解；其中的解类对照帮助回查数据、解空间和方程意义的区别。伴读时，Walter A. Strauss 的《Partial Differential Equations: An Introduction》\cite[Sections 2.1--2.2, 3.2, 3.4]{Strauss2008Partial}分别讨论一维波动、因果性与能量、半直线反射及外力问题；Fritz John 的《Partial Differential Equations》\cite{John1991Partial}中关于特征线和高维波方程的论述，可接在本章之后阅读。

% !TeX root = ../../Book2.tex
\section{齐次方程}
\label{b2:sec:0901}

一根均匀弦作小幅横向振动时，位移的时间加速度与空间曲率成正比。忽略阻尼和外力，得到 \(u_{tt}-c^2u_{xx}=0\)，其中 \(c>0\) 为固定波速。本章先在整条实线上讨论，因而没有端点反射。方程含二阶时间导数，初始状态必须同时指定形状与速度；只给初始形状不会确定随后的运动。

\subsection{特征变量}
\label{b2:subsec:090101}

常系数使波算子可以分解为两个可交换的一阶算子：
\[
 \partial_t^2-c^2\partial_x^2
 =(\partial_t-c\partial_x)(\partial_t+c\partial_x).
\]
若令 \(\alpha=x+ct\)、\(\beta=x-ct\)，则
\[
 x=\dfrac{\alpha+\beta}{2},\quad
 t=\dfrac{\alpha-\beta}{2c},\quad
 \partial_t^2-c^2\partial_x^2=-4c^2\partial_\alpha\partial_\beta.
\]
这对坐标称为\term[tezheng-bianliang]{特征变量}[characteristic variables]；两族坐标线正是斜率分别为 \(\pm c\) 的传播直线。在一个坐标矩形上，方程成为 \(U_{\alpha\beta}=0\)。先对 \(\beta\) 积分，再对 \(\alpha\) 积分，得到 \(U=A(\alpha)+B(\beta)\)。相应地，原坐标中的两个分量分别保持形状向左、向右移动。

\ProseParagraphBreak

对初值问题，直接追踪一阶导数更便于把解与数据联系起来。设 \(u\) 为二次连续可微的解，记
\[
 p=u_t+cu_x,\quad q=u_t-cu_x.
\]
由原方程，\(p_t-cp_x=0\)、\(q_t+cq_x=0\)。沿通过 \((t,x)\) 的直线积分，便有
\[
 p(t,x)=p(0,x+ct),\quad q(t,x)=q(0,x-ct).
\]
这里的符号可由沿线求导检查：曲线 \(s\mapsto(s,x+c(t-s))\) 的导数算子是 \(\partial_t-c\partial_x\)。因此 \(p\) 沿向左的直线传递，\(q\) 沿向右的直线传递。两者不是额外的数据，而是位移梯度和速度的线性组合。

\subsection{d'Alembert 公式}
\label{b2:subsec:090102}

考虑初值
\[
 u(0,x)=f(x),\quad u_t(0,x)=g(x).
\]
初始组合为 \(p(0,x)=g(x)+cf'(x)\)、\(q(0,x)=g(x)-cf'(x)\)。将它们相加，恢复 \(u_t\)，再从 \(0\) 积分到 \(t\)，就会出现区间 \([x-ct,x+ct]\) 上的初速度积分。

\begin{theorem}[d'Alembert]\label{b2:thm:09-dalembert}
若 \(f\in C^2(\mathbb R)\)、\(g\in C^1(\mathbb R)\)，则初值问题在 \(\mathbb R_t\times\mathbb R_x\) 上有唯一经典解，且
\begin{equation}\label{b2:eq:09-dalembert}
 u(t,x)=\dfrac{f(x+ct)+f(x-ct)}2
 +\dfrac1{2c}\int_{x-ct}^{x+ct}g(y)\,\mathrm{d}y.
\end{equation}
当 \(t<0\) 时，积分按有向区间解释；不要求 \(f\)、\(g\) 在无穷远有界或衰减。
\end{theorem}
\begin{proof}
由前述沿线等式，任一经典解都满足
\[
 u_t(t,x)=\dfrac{g(x+ct)+g(x-ct)}2
 +\dfrac c2\bigl(f'(x+ct)-f'(x-ct)\bigr).
\]
把右端对时间积分并加上 \(f(x)\)。两个导数项分别换元，得到两个平移的平均减去 \(f(x)\)；速度项换元得到\eqref{b2:eq:09-dalembert}的区间积分。因此任何解都必须是该表达式，唯一性已经得到。

反过来，由 \(f\in C^2\)、\(g\in C^1\)，公式右端二次连续可微。位移项的二阶时间导数等于其二阶空间导数乘 \(c^2\)。若将积分项记作 \(v\)，则
\[
 v_t=\dfrac{g(x+ct)+g(x-ct)}2,\quad
 v_x=\dfrac{g(x+ct)-g(x-ct)}{2c},
\]
从而 \(v_{tt}=c^2v_{xx}\)。令 \(t=0\)，积分项与位移项的时间导数分别给出 \(g\) 与 \(0\)，位移本身则为 \(f\)。所以存在性和全部初值条件成立。
\end{proof}

表达式\eqref{b2:eq:09-dalembert}称为\term[dalembert-gongshi]{d'Alembert 公式}[d'Alembert's formula]。它把初始位移分成两个等幅平移，把初始速度转化为一个不断扩张的区间积分。初始位移只在区间的两个端点取值，初始速度却在整段区间内参与；这一区别将影响长时间留下的位移。

\ProseParagraphBreak

例如取 \(g=0\)、\(f(x)=e^{-x^2}\)，初始隆起分裂成两个各有原来一半高度的平移。若取 \(f(x)=\cos(kx)\)、\(g=0\)，则
\[
 u(t,x)=\cos(kx)\cos(ckt).
\]
后一种运动的零点位置不变，通常称为\term[zhubo]{驻波}[standing wave]；前一种保持轮廓平移的分量称为\term[xingbo]{行波}[traveling wave]。相同方程容许不同运动，区别来自初值，而不来自方程阶数。

\subsection{初值问题}
\label{b2:subsec:090103}

为了只得到向右的行波 \(u(t,x)=f(x-ct)\)，必须取 \(g=-cf'\)。若只把弦拉成 \(f\) 然后静止释放，即令 \(g=0\)，通常会同时产生两个方向的波。一般数据可用一个原函数 \(G'=g\) 写成
\[
 u(t,x)=A(x+ct)+B(x-ct),\quad
 A=\dfrac f2+\dfrac G{2c},\quad
 B=\dfrac f2-\dfrac G{2c}.
\]
原函数的常数分别改变 \(A\)、\(B\) 的常数部分，但不改变它们的和。这说明行波分解具有一个无关紧要的常数自由度，初值问题的解本身仍唯一。

\ProseParagraphBreak

初值线 \(t=0\) 同时与两族特征线相交，因而能确定两个传播分量。若改在特征线 \(x=ct\) 上给值，则沿这条线只能看到 \(A(2ct)+B(0)\)，不能恢复整个函数 \(B\)。例如 \(B(x-ct)\) 在该线上恒为零，只需 \(B(0)=0\)，但离开该线以后可以非零。对于二阶方程，数据的几何位置与数据数量同样重要。

\ProseParagraphBreak

公式还表明齐次运动在时间上可逆：把 \(t\) 换成 \(-t\)，得到的数据为 \((f,-g)\)。这种可逆性与后面热方程的平滑过程不同。若初始轮廓有折角，把它平移并不会使折角消失；在经典可微性不足时，仍可在分布或能量意义下解释同一公式。\secref{b2:sec:0903}将给出这一延拓的空间、极限和唯一性，而不是仅把不可微函数代入经典方程。

% !TeX root = ../../Book2.tex
\section{非齐次方程}
\label{b2:sec:0902}

外力在每一时刻都改变速度。线性叠加允许把持续作用拆成各个时刻产生的初速度增量，再分别传播到观察时刻。积分上限处的贡献决定了非齐次项；若只对积分内部求导，便会漏掉外力本身。

\subsection{Duhamel 原理}
\label{b2:subsec:090201}

先把零位移、给定速度的解算子记为
\[
 S_c(t)h(x)\coloneqq\dfrac1{2c}\int_{x-ct}^{x+ct}h(y)\,\mathrm{d}y.
\]
它满足 \(S_c(0)h=0\)、\(\partial_tS_c(0)h=h\)。若时刻 \(s\) 的力为 \(F(s,\cdot)\)，它在时长 \(t-s\) 内产生的响应是 \(S_c(t-s)F(s,\cdot)\)。对这些响应积分的做法称为\term[duhamel-yuanli]{Duhamel 原理}[Duhamel's principle]。
\symidx{\(S_c(t)\)：一维波方程的初速度传播算子}

\begin{theorem}[Duhamel，一维波方程]\label{b2:thm:09-duhamel}
设 \(T>0\)，\(f\in C^2(\mathbb R)\)、\(g\in C^1(\mathbb R)\)、\(F\in C^1([0,T]\times\mathbb R)\)。问题
\[
 u_{tt}-c^2u_{xx}=F,\quad u(0)=f,\quad u_t(0)=g
\]
有唯一经典解；它等于\eqref{b2:eq:09-dalembert}的齐次解加上
\begin{equation}\label{b2:eq:09-duhamel}
 w(t,x)=\int_0^tS_c(t-s)F(s,\cdot)(x)\,\mathrm{d}s
 =\dfrac1{2c}\int_0^t\int_{x-c(t-s)}^{x+c(t-s)}F(s,y)\,\mathrm{d}y\,\mathrm{d}s.
\end{equation}
\end{theorem}
\begin{proof}
记 \(V(\tau,s,x)=S_c(\tau)F(s,\cdot)(x)\)。固定 \(s\) 时，\(V_{\tau\tau}=c^2V_{xx}\)，且 \(V(0,s,x)=0\)、\(V_\tau(0,s,x)=F(s,x)\)。各阶求导只涉及一个紧三角形上的 \(F\)、\(F_x\)，故可在积分号下求导。Leibniz 公式给出
\[
 w_t=\int_0^tV_\tau(t-s,s,x)\,\mathrm{d}s,\quad
 w_{tt}=F(t,x)+\int_0^tV_{\tau\tau}(t-s,s,x)\,\mathrm{d}s.
\]
另一方面，\(c^2w_{xx}\) 正是第二式的积分项。所以 \(w_{tt}-c^2w_{xx}=F\)。积分区间在 \(t=0\) 退化，给出 \(w(0)=w_t(0)=0\)。减去这个特解后使用\cref{b2:thm:09-dalembert}，即可得到一般初值的存在性与唯一性。
\end{proof}

外力项的积分区域是观察点向过去展开的三角形。对固定 \(s\)，水平截面宽度为 \(2c(t-s)\)：较早施加的力能够从较远的位置传来，刚施加的力只能从观察点附近传来。这一几何结构比公式中的常数更有普遍性，后面的高维问题将把区间改成球。

\subsection{外力项}
\label{b2:subsec:090202}

若 \(F(s,y)=a(s)\) 与空间无关，且初值为零，公式化为
\[
 u(t,x)=\int_0^t(t-s)a(s)\,\mathrm{d}s.
\]
因此恒力产生 \(u=at^2/2\)，恰与没有空间变化时的常微分方程 \(u_{tt}=a\) 一致。常数 \(1/(2c)\) 被区间长度中的 \(2c\) 抵消，这也是检查归一化的一种简便方法。

\ProseParagraphBreak

外力的大小受到控制，并不意味着位移在所有时间上一致有界。取 \(k>0\)、\(\omega=ck\)，令
\[
 F(t,x)=\cos(\omega t)\cos(kx),\quad f=g=0.
\]
设 \(u(t,x)=a(t)\cos(kx)\)，则 \(a''+\omega^2a=\cos(\omega t)\)，从而
\[
 a(t)=\dfrac{t\sin(\omega t)}{2\omega}.
\]
直接求导可核对方程与零初值。外力的时间频率与自由振动频率相同，持续作用可使振幅增长；这种现象称为\term[gongzhen]{共振}[resonance]。本例为空间周期的经典解，不具有整条实线上的有限总能量，因此不把它代入稍后的全空间能量积分。

\ProseParagraphBreak

若仅需有限时间的逐点控制，\eqref{b2:eq:09-duhamel}直接给出
\[
 \|w(t)\|_\infty\leq\int_0^t(t-s)\|F(s)\|_\infty\,\mathrm{d}s.
\]
在 \(F\) 一致有界时右端不超过 \(t^2\|F\|_\infty/2\)。该估计没有利用振荡抵消，适合控制局部存在时间内的误差；它不声称捕捉每一个振荡外力的精确长时间增长。

\subsection{有限传播速度}
\label{b2:subsec:090203}

给定 \(t_0>0\)、\(x_0\in\mathbb R\)，记过去的传播区域为
\[
 K_{t_0,x_0}\coloneqq
 \{\,(s,y):0\leq s\leq t_0,\ |y-x_0|\leq c(t_0-s)\,\}.
\]
若称某个区域为观察点的\term[yilai-quyu]{依赖区域}[domain of dependence]，意思是该处的解只由区域内的数据确定。相反，从一个初始位置出发可能受到影响的未来区域称为\term[yingxiang-quyu]{影响区域}[domain of influence]。这两个说法分别固定终点和起点，不能把箭头方向混同。

\begin{proposition}[有限传播速度]\label{b2:prop:09-finite-speed}
设两个经典解具有相同波速。若它们的初始位移、初始速度在 \([x_0-ct_0,x_0+ct_0]\) 上分别相等，外力在 \(K_{t_0,x_0}\) 上相等，则两解在整个 \(K_{t_0,x_0}\) 上相等。特别，若初值及每个时刻的外力都支撑在固定闭集 \(A\) 内，则
\[
 \operatorname{supp}u(t,\cdot)\subseteq
 \{\,x:\operatorname{dist}(x,A)\leq ct\,\},\quad 0\leq t\leq T.
\]
\end{proposition}
\begin{proof}
对两解作差。区域中任一点 \((s,y)\) 的初始依赖区间包含在 \([x_0-ct_0,x_0+ct_0]\) 内，而它的过去三角形包含在 \(K_{t_0,x_0}\) 内：这两点都由三角不等式
\(|z-x_0|\leq|z-y|+|y-x_0|\) 得出。因此 d'Alembert 公式中的两端点值与区间积分为零，Duhamel 积分也为零。

若 \(\operatorname{dist}(x,A)>ct\)，则过去三角形的任何截面都不与 \(A\) 相交，故解在该点为零。距离函数连续，使这个零值结论在点的邻域也成立；于是得到所陈述的支集包含关系。
\end{proof}

\cref{b2:fig:09-cone}画出了同一个区域的两种用途：计算 \((t_0,x_0)\) 的位移时读取过去的数据；证明唯一性时则比较一整个收缩区间上的能量。第二种方式不依赖显式公式，因而能够推广到其他维数。
\begin{figure}[htbp]
\centering
\begin{tikzpicture}[x=1cm,y=1cm,>=stealth]
\fill[AnalysisBlue!9] (-3,0)--(0,3)--(3,0)--cycle;
\draw[->] (-3.8,0)--(3.8,0) node[right]{\(x\)};
\draw[->] (-3.6,-.2)--(-3.6,3.5) node[above]{\(t\)};
\draw[thick,AnalysisBlue] (-3,0)--(0,3)--(3,0);
\draw[dashed] (-1.8,1.2)--(1.8,1.2);
\fill (0,3) circle(1.7pt) node[above right]{\((t_0,x_0)\)};
\node[below] at (-3,0){\(x_0-ct_0\)};
\node[below] at (3,0){\(x_0+ct_0\)};
\node[below] at (0,0){\(x_0\)};
\node at (0,.6){\(K_{t_0,x_0}\)};
\node[above] at (0,1.2){时刻 \(s\) 的截面};
\end{tikzpicture}
\caption[一维波动方程的依赖三角形]{一维波动方程的依赖三角形。截面半径为 \(c(t_0-s)\)，两侧边界随时间以速度 \(c\) 向观察点收缩；图中横纵比例仅为示意。}
\label{b2:fig:09-cone}
\end{figure}


传播速度的结论只给出支集的外界，不保证界内处处非零。左右行波可以抵消，数据本身也可能具有空隙。此外，在一维即使初始数据紧支撑，固定位置也未必在波前经过以后回到零位移：若 \(f=0\)、\(g\) 紧支撑，则对足够大的 \(t\)，
\[
 u(t,x)=\dfrac1{2c}\int_{\mathbb R}g(y)\,\mathrm{d}y.
\]
这个尾部不违反有限传播，因为它位于允许的传播区域内部。三维方程在相同条件下具有更强的无尾传播性质，须到\chapref{b2:ch:11}重新证明。

% !TeX root = ../../Book2.tex
\section{能量方法}
\label{b2:sec:0903}

显式公式给出解的每个点值；能量方法则把速度与空间梯度放进一个积分量。后者能在点值不可定义、二阶导数只存在于分布意义下时继续工作。真正需要保存的是估计所控制的范数，以及把光滑解送到极限的过程。

\subsection{能量恒等式}
\label{b2:subsec:090301}

先设解和外力光滑，且在每个有界时间区间内有共同的紧空间支集。定义\term[nengliang-midu]{能量密度}[energy density]及总能量
\[
 e(t,x)\coloneqq\dfrac12\bigl(|u_t(t,x)|^2+c^2|u_x(t,x)|^2\bigr),\quad
 E(t)\coloneqq\int_{\mathbb R}e(t,x)\,\mathrm{d}x.
\]
复值解也使用这一实数能量。把方程乘以 \(\overline{u_t}\) 并取实部，得到局部恒等式
\begin{equation}\label{b2:eq:09-local-energy}
 \partial_te-\partial_x\bigl(c^2\operatorname{Re}(u_x\overline{u_t})\bigr)
 =\operatorname{Re}(F\overline{u_t}).
\end{equation}
因此相应的\term[nengliang-tongliang]{能量通量}[energy flux]为 \(-c^2\operatorname{Re}(u_x\overline{u_t})\)。负号与守恒式 \(\partial_te+\partial_xj=\operatorname{Re}(F\overline{u_t})\) 相配。
\symidx{\(E(t)\)：波动方程的总能量}

\begin{proposition}[能量与外力功]\label{b2:prop:09-energy}
在上述光滑紧支集假设下，
\begin{equation}\label{b2:eq:09-energy-identity}
 E(t)=E(0)+\operatorname{Re}\int_0^t\int_{\mathbb R}
 F(s,x)\overline{u_t(s,x)}\,\mathrm{d}x\,\mathrm{d}s.
\end{equation}
特别，\(F=0\) 时总能量守恒。一般有
\begin{equation}\label{b2:eq:09-energy-bound}
 \sqrt{2E(t)}\leq\sqrt{2E(0)}+\int_0^t\|F(s)\|_2\,\mathrm{d}s.
\end{equation}
\end{proposition}
\begin{proof}
将\eqref{b2:eq:09-local-energy}对空间积分，支集条件消去无穷远边界项，再对时间积分，得到恒等式。由 Cauchy--Schwarz，
\[
 E'(t)\leq\|F(t)\|_2\cdot\|u_t(t)\|_2
 \leq\sqrt{2E(t)}\cdot\|F(t)\|_2.
\]
不能在 \(E(t)=0\) 时直接除以 \(\sqrt{E(t)}\)。对 \(\varepsilon>0\)，改对 \(\sqrt{2E(t)+\varepsilon^2}\) 求导，其导数不超过 \(\|F(t)\|_2\)。积分后令 \(\varepsilon\downarrow0\)，得到估计。
\end{proof}

齐次解的动能 \(\|u_t\|_2^2/2\) 与势能 \(c^2\|u_x\|_2^2/2\) 可以互相转换，守恒的是两者之和。总位移的 \(L^2\) 范数并不出现在能量里，因此不能直接从 \(E(t)=E(0)\) 得到 \(\|u(t)\|_2=\|f\|_2\)。其零阶部分须另由时间积分控制。

\subsection{唯一性}
\label{b2:subsec:090302}

全空间积分要求有限总能量，但经典解唯一性不应排除多项式增长的初值。把能量放在随时间收缩的有限区间上，可以摆脱无穷远条件。固定 \(x_0\)、\(t_0\)，设
\[
 a(t)=x_0-c(t_0-t),\quad b(t)=x_0+c(t_0-t),\quad
 E_K(t)=\int_{a(t)}^{b(t)}e(t,x)\,\mathrm{d}x.
\]
移动端点的速度为 \(a'=c\)、\(b'=-c\)，因此即使 \(F=0\)，也不能仅把空间导数积分后宣称局部能量不变。

\begin{proposition}[收缩区间的能量]\label{b2:prop:09-cone-energy}
若 \(u\) 在依赖三角形的邻域中为经典齐次解，则对 \(0<t<t_0\)，
\[
 E_K'(t)=-\dfrac c2|u_t-cu_x|^2(t,b(t))
             -\dfrac c2|u_t+cu_x|^2(t,a(t))\leq0.
\]
若初始位移与初始速度在底边区间上为零，则 \(u\) 在整个三角形上为零。
\end{proposition}
\begin{proof}
Leibniz 公式与\eqref{b2:eq:09-local-energy}给出
\[
 E_K'=c^2\operatorname{Re}(u_x\overline{u_t})(t,b)-c\cdot e(t,b)
       -c^2\operatorname{Re}(u_x\overline{u_t})(t,a)-c\cdot e(t,a).
\]
分别配方即得所示平方和。若初值为零，则底边内部的 \(f'\) 也为零，故 \(E_K(0)=0\)。非负性与单调性给出 \(E_K(t)=0\)，从而内部的 \(u_t\)、\(u_x\) 都为零。对内部任一点沿竖直线积分到初值线，这条线始终留在三角形中，故 \(u(t,x)=u(0,x)=0\)。边界和顶点由连续性得到。
\end{proof}

把命题用于两个解之差，可重新得到局部唯一性和有限传播速度。这一次不使用 d'Alembert 表达式，核心是侧边通量恰被边界收缩带走的能量控制。选择比 \(c\) 更慢的收缩速度时，平方配方的符号可能失效；波速是估计所允许的几何尺度。

\subsection{稳定性}
\label{b2:subsec:090303}

比较两个光滑解 \(u\)、\(\widetilde u\)，分别记数据差为 \(h=f-\widetilde f\)、\(k=g-\widetilde g\)、\(G=F-\widetilde F\)。对差解使用能量估计，令
\[
 M_T\coloneqq\bigl(\|k\|_2^2+c^2\|h'\|_2^2\bigr)^{1/2}
                +\int_0^T\|G(s)\|_2\,\mathrm{d}s.
\]
则在 \(0\leq t\leq T\) 上有
\begin{equation}\label{b2:eq:09-stability}
 \|u_t(t)-\widetilde u_t(t)\|_2\leq M_T,\quad
 c\|u_x(t)-\widetilde u_x(t)\|_2\leq M_T,\quad
 \|u(t)-\widetilde u(t)\|_2\leq\|h\|_2+tM_T.
\end{equation}
最后一项来自 \(u(t)-\widetilde u(t)=h+\int_0^t(u_s-\widetilde u_s)\,\mathrm{d}s\)。它补上了能量没有控制的零阶部分，也解释了估计为何依赖所观察的有限时间长度。

\ProseParagraphBreak

以上结果给出的\term[lianxu-yilai]{连续依赖}[continuous dependence]有明确含义：初始位移在 \(H^1\) 范数、初始速度在 \(L^2\) 范数、外力在时间可积的 \(L^2\) 范数中接近，则位移与速度在整个有限时间区间内一致接近。这里不能仅保留 \(\|h\|_2\)。例如取固定非零光滑紧支集函数 \(\eta\)，令 \(h_n(x)=n^{-1}\eta(x)\sin(nx)\)。虽然 \(h_n\to0\) 于 \(L^2\)，其导数含 \(\eta(x)\cos(nx)\)，平方范数趋于 \(\|\eta\|_2^2/2\)；后一极限由 Riemann--Lebesgue 定理处理振荡项得到。很小的位移可以携带不消失的高频能量。

\subsection{低正则初值}
\label{b2:subsec:090304}

以下使用第一编已经建立的 \(L^2\) 与弱导数语言。记 \(H^1(\mathbb R)=\{f\in L^2:f'\in L^2\}\)，范数平方为 \(\|f\|_2^2+\|f'\|_2^2\)。外力空间 \(L^1(0,T;L^2)\) 指强可测且 \(\int_0^T\|F(s)\|_2\,\mathrm{d}s<\infty\) 的函数；其中的向量值积分可由简单函数在这一积分范数中的逼近定义。积分三角不等式保证定义不依赖逼近。

\begin{definition}\label{b2:def:09-energy-solution}
设 \(f\in H^1\)、\(g\in L^2\)、\(F\in L^1(0,T;L^2)\)。若
\[
 u\in C([0,T];H^1)\cap C^1([0,T];L^2),
\]
在 \((0,T)\times\mathbb R\) 上满足分布方程 \(u_{tt}-c^2u_{xx}=F\)，并在所示范数意义下满足 \(u(0)=f\)、\(u_t(0)=g\)，则称 \(u\) 为本初值问题的\term[nengliang-jie]{能量解}[energy solution]。
\end{definition}

分布方程是双线性测试等式
\[
 \int_0^T\int_{\mathbb R}u(\varphi_{tt}-c^2\varphi_{xx})\,\mathrm{d}x\,\mathrm{d}t
 =\int_0^T\int_{\mathbb R}F\varphi\,\mathrm{d}x\,\mathrm{d}t,
 \quad \varphi\in C_c^\infty((0,T)\times\mathbb R).
\]
这里不放复共轭；能量计算中的共轭来自 Hilbert 内积，两种配对各有用途。初值通过时间连续性另行给定，不能由仅在开时间区间内成立的测试等式推出来。

\begin{theorem}[能量极限]\label{b2:thm:09-energy-solution}
上述数据存在唯一能量解。光滑紧支集数据在 \(H^1\times L^2\times L^1(0,T;L^2)\) 中趋于给定数据时，对应经典解在 \(C([0,T];H^1)\cap C^1([0,T];L^2)\) 中趋于该解。能量恒等式\eqref{b2:eq:09-energy-identity}及稳定性估计\eqref{b2:eq:09-stability}仍成立。
\end{theorem}
\begin{proof}
先说明可选取所需数据。对 \(f\) 乘以在 \([-R,R]\) 上等于一、导数为 \(O(R^{-1})\) 的光滑截断，再卷积磨光函数。尾部的 \(f\)、\(f'\) 平方积分趋零，截断导数乘 \(f\) 的范数也趋零；卷积由平移连续性在 \(L^2\) 中同时逼近函数和导数。于是得到 \(f_n\in C_c^\infty\) 在 \(H^1\) 中逼近 \(f\)。对 \(g\) 只需同样的零阶论证。外力先用有限个 \(L^2\) 值简单函数逼近，再逼近这些值的空间函数和各个时间示性函数，得到光滑且空间紧支撑的 \(F_n\)，误差趋零于 \(L^1(0,T;L^2)\)。

由 d'Alembert--Duhamel 公式产生经典解 \(u_n\)。稳定性估计说明 \(u_n\)、\(\partial_xu_n\)、\(\partial_tu_n\) 分别在 \(C([0,T];L^2)\) 中为 Cauchy 列，记极限为 \(u,w,v\)。空间分部积分取极限得到 \(w=\partial_xu\)；等式 \(u_n(t)=f_n+\int_0^t\partial_su_n(s)\,\mathrm{d}s\) 在 \(L^2\) 中取极限得到 \(u(t)=f+\int_0^tv(s)\,\mathrm{d}s\)，故 \(u_t=v\)。测试等式的各项由范数收敛取极限，得到分布方程。初值也由一致时间收敛保留。

能量中的两个平方范数随强收敛趋于极限。功积分的差可分成 \(F_n-F\) 乘 \(\overline{u_{n,t}}\) 和 \(F\) 乘 \(\overline{u_{n,t}-u_t}\)；分别由外力的积分范数收敛和速度的一致范数收敛控制。因此能量恒等式传到极限。相同过程也保留稳定性估计，并说明极限不依赖所选数据逼近。

还须排除其他分布能量解。设 \(z\) 为零数据的能量解，对空间卷积光滑紧支集核，得到 \(z_\varepsilon\)。它在空间具有任意阶 \(L^2\) 导数，而且 \(z_{\varepsilon,tt}=c^2z_{\varepsilon,xx}\)。右端连续取值于 \(L^2\)，故 \(z_{\varepsilon,t}\) 在时间上可按 \(L^2\) 微积分求导。Hilbert 内积的乘积法则与全空间分部积分给出其能量守恒；全空间分部积分可先乘空间截断，误差由 \(L^2\) 尾部与 Cauchy--Schwarz 控制后趋零。零初值使能量为零，因而 \(z_{\varepsilon,t}=0\)，再由 \(z_\varepsilon(0)=0\) 得到 \(z_\varepsilon=0\)。让 \(\varepsilon\downarrow0\)，得到 \(z=0\)。
\end{proof}

有限传播也传到能量极限，但支集应按 \(L^2\) 函数的本性支集理解。若初值和几乎每个时刻的外力支撑在同一闭集 \(A\) 内，可以先截断再用半径 \(\varepsilon\) 的磨光核逼近，使近似数据支撑在 \(A\) 的 \(\varepsilon\) 邻域。经典解支撑在距离 \(A\) 不超过 \(ct+\varepsilon\) 的区域内。在其外部固定一个紧集，先令逼近指标趋于无穷，再令 \(\varepsilon\) 趋零，\(L^2\) 极限仍为零。

\ProseParagraphBreak

最后取 \(f(x)=\max\{1-|x|,0\}\)、\(g=0\)。它的弱导数分段为常数并属于 \(L^2\)，所以数据属于能量空间；公式仍为两个帐篷形状的平移平均。初始的三个折角分别沿特征线前进，得到连续但通常非二次可微的解。能量解的存在并不制造额外的平滑性，它给原先不满足经典正则性的数据提供了稳定且唯一的演化意义。


\begin{chaptersummary}
常系数波算子分成两个一阶传播算子，使初始速度与位移梯度沿两族直线移动。d'Alembert 公式由此确定齐次演化；Duhamel 原理把时空外力积累为过去三角形上的积分。公式依赖的区域有限，而区域内部是否存在尾部还取决于数据和维数。

\ProseParagraphBreak

能量方法把同一几何转化为通量不等式。全空间能量控制速度与梯度，位移的零阶范数另由时间积分补齐；收缩区间上的能量则给出无需无穷远假设的局部唯一性。共同紧空间支集可直接消去无穷远通量，对一般有限能量解则须先用空间截断证明误差趋零。完整 \(H^1\times L^2\) 状态的稳定性与仅观察位移的 \(L^2\) 稳定性，控制的是不同层次的误差。

\ProseParagraphBreak

低正则解的构造依次完成三个任务：稳定估计使光滑数据的解收敛，强收敛把方程与能量恒等式传到极限，空间磨光再证明任意能量解的唯一性。有限区间的特征函数展开和高维波方程，也将用稳定估计构造低正则解。
\end{chaptersummary}

\BookExercises

前几题比较传播方向、初速度和端点反射，随后讨论守恒量、能量分配与外力冲量。最后两题分别改变系数和所用范数，用以检验能量方法的适用范围。

% 自拟：两族传播量的能量分配和单向条件，不只重写d'Alembert公式。
\Needspace{6\baselineskip}
\begin{exercise}\label{b2:ex:09-directional-energy}
设 \(c>0\)、\(\mathbb K\in\{\mathbb R,\mathbb C\}\)，\(f\in H^1(\mathbb R;\mathbb K)\)、\(g\in L^2(\mathbb R;\mathbb K)\)。令 \(u\) 为 \(u_{tt}-c^2u_{xx}=0\)、\(u(0)=f\)、\(u_t(0)=g\) 在 \(t\geq0\) 上的能量解，即其在每个有限时间区间都属于 \(C_tH^1_x\cap C_t^1L^2_x\)。记 \(E(t)=\frac12(\|u_t(t)\|_2^2+c^2\|u_x(t)\|_2^2)\)。证明
\[
 E(t)=\dfrac14\|g+cf'\|_2^2+\dfrac14\|g-cf'\|_2^2.
\]
分别解释两项的传播方向，并证明对所有 \(t\geq0\) 有 \(u(t)=f(\cdot-ct)\) 当且仅当 \(g=-cf'\) 于 \(L^2(\mathbb R)\)。
\end{exercise}
\begin{solution}
光滑情形中 \(p=u_t+cu_x\)、\(q=u_t-cu_x\) 分别是 \(g+cf'\)、\(g-cf'\) 的左、右平移。恒等式 \(|u_t|^2+c^2|u_x|^2=(|p|^2+|q|^2)/2\) 与平移等距性给出所求能量分解。对光滑数据逼近取极限，它也适用于能量解。若 \(g=-cf'\)，则左行的 \(p\) 恒为零，右平移 \(f(x-ct)\) 在能量空间中具有相同初值，由唯一性得到等式。反过来，右平移在 \(L^2\) 中的时间导数在零时刻为 \(-cf'\)，故初速度只能取该值。
\end{solution}

% 自拟：阶梯初速度产生连续梯形响应，核对支集和一维尾部。
\Needspace{6\baselineskip}
\begin{exercise}\label{b2:ex:09-rectangular-velocity}
设实数 \(c>0\)、\(a>0\)，\(f=0\)、\(g=\mathbf1_{[-a,a]}\)。对 \(\mathbb R_x\) 上满足 \(u_{tt}-c^2u_{xx}=0\)、\(u(0)=f\)、\(u_t(0)=g\) 的实值能量解，求 \(t\geq0\) 时位移的连续代表，描述其空间支集、最大值及固定空间点的长时间极限，并计算总能量 \(E(t)=\frac12(\|u_t(t)\|_2^2+c^2\|u_x(t)\|_2^2)\)。
\end{exercise}
\begin{solution}
用初速度积分可写为
\[
 u(t,x)=\dfrac1{2c}\max\{0,\min\{x+ct,a\}-\max\{x-ct,-a\}\},\quad t\geq0.
\]
这是两个区间的交长除以 \(2c\)，故连续；当 \(t>0\) 时，它为正恰在 \(|x|<a+ct\) 内，空间支集为闭区间 \([-a-ct,a+ct]\)；\(t=0\) 时位移恒零，支集为空。交长的最大值为 \(2\min\{a,ct\}\)，所以最大值为 \(\min\{a/c,t\}\)。固定 \(x\) 后，当 \(ct>a+|x|\) 时整个初速度区间已被包含，故解恒为 \(a/c\)。初始能量为 \(\|g\|_2^2/2=a\)，能量守恒给出以后相同的值。正位移的平台扩大时，能量仍集中在移动的过渡区域。
\end{solution}

% 自拟：把全线解用于半线Dirichlet反射，增加相容性检验。
\Needspace{6\baselineskip}
\begin{exercise}\label{b2:ex:09-halfline-reflection}
设 \(c>0\)、\(\mathbb K\in\{\mathbb R,\mathbb C\}\)，\(f\in C_c^2([0,\infty);\mathbb K)\)、\(g\in C_c^1([0,\infty);\mathbb K)\)。这里半线端点的导数取右导数，支集只要求在半线中紧。另设 \(f(0)=f''(0)=g(0)=0\)。通过奇延拓求解 \(t,x\geq0\) 上的 \(C^2\) 半线问题
\[
 u_{tt}-c^2u_{xx}=0,\quad u(t,0)=0,\quad
 u(0,x)=f(x),\quad u_t(0,x)=g(x).
\]
写出 \(0\leq x<ct\) 时的公式，并说明反射改变了哪个行波的符号。
\end{exercise}
\begin{solution}
令 \(f_o(x)=f(x)\) 于 \(x\geq0\)，\(f_o(x)=-f(-x)\) 于 \(x<0\)，对 \(g\) 类似定义。所给相容条件使 \(f_o\in C^2\)、\(g_o\in C^1\)：奇延拓的二阶导数两侧极限分别为 \(\pm f''(0)\)。全线解保持对空间的奇性，所以在原点为零。若 \(x<ct\)，分开速度积分的正负两段，得到
\[
 u(t,x)=\dfrac{f(x+ct)-f(ct-x)}2
 +\dfrac1{2c}\int_{ct-x}^{ct+x}g(y)\,\mathrm{d}y.
\]
尚未碰到端点时使用普通d'Alembert公式；碰到端点以后，原本继续走向负半线的位移分量由镜像返回，并带相反符号。因此固定端点的反射反转位移符号。若两个半线经典解有相同数据，在截断的依赖区域使用能量恒等式，原点处 \(u_t=0\) 消去通量，得到同样的唯一性。
\end{solution}

% 自拟：守恒动量与能量的最优不等式，等号刻画单向传播。
\Needspace{6\baselineskip}
\begin{exercise}\label{b2:ex:09-momentum}
设 \(c>0\)，\(u\in C^2([0,\infty)\times\mathbb R;\mathbb R)\) 满足 \(u_{tt}-c^2u_{xx}=0\)。假设对每个 \(T>0\)，存在 \(R_T>0\) 使 \(u(t,x)=0\) 于 \(0\leq t\leq T\)、\(|x|>R_T\)；这是有限时间带上共同的紧空间支集条件，不要求整个时空支集紧。定义
\[
 P(t)=-\int_{\mathbb R}u_tu_x\,\mathrm{d}x,\qquad
 E(t)=\frac12\int_{\mathbb R}(u_t^2+c^2u_x^2)\,\mathrm{d}x.
\]
证明 \(P\) 守恒、\(c|P|\leq E\)，并刻画 \(cP=E\) 的情形；用非零解说明不等式中的常数最优。
\end{exercise}
\begin{solution}
分部积分给出
\[
 P'=-\int(c^2u_{xx}u_x+u_tu_{xt})\,\mathrm{d}x
 =-\dfrac12\int\partial_x(c^2u_x^2+u_t^2)\,\mathrm{d}x=0.
\]
共同的空间支集保证求导及上述分部积分均在一个固定紧区间中进行。由 \(2c|u_tu_x|\leq u_t^2+c^2u_x^2\)，积分得到不等式。更精确地，\(E-cP=\frac12\int(u_t+cu_x)^2\)。总能量与 \(P\) 都守恒，故一个时刻满足 \(cP=E\)，就会在所有时刻满足；非负连续被积函数的积分为零，又给出 \(u_t+cu_x=0\) 处处成立。沿向右的直线积分，得到 \(u(t,x)=f(x-ct)\)，其中 \(f=u(0,\cdot)\in C_c^2(\mathbb R;\mathbb R)\)。反过来，这样的右行波确实取等。取任意非零 \(f\in C_c^\infty(\mathbb R;\mathbb R)\)，则 \(f'\) 不恒为零，因而 \(E>0\)，等号说明常数不能改善。这个例子的支集随时间平移，但在任意有限时间带上仍有共同的紧空间支集。
\end{solution}

% 自拟：紧支集数据的有限时刻能量均分，区别于位移消失。
\Needspace{6\baselineskip}
\begin{exercise}\label{b2:ex:09-equipartition}
设 \(c>0\)、\(R>0\)，\(f\in C_c^2(\mathbb R;\mathbb R)\)、\(g\in C_c^1(\mathbb R;\mathbb R)\) 均支撑在 \([-R,R]\) 内，\(u\) 是全实线齐次波动方程 \(u_{tt}-c^2u_{xx}=0\)、\(u(0)=f\)、\(u_t(0)=g\) 的经典解。记动能为 \(K(t)=\|u_t(t)\|_2^2/2\)，势能为 \(V(t)=c^2\|u_x(t)\|_2^2/2\)，总能量为 \(E=K+V\)。证明当 \(t>R/c\) 时，\(K(t)=V(t)=E/2\)。即使 \(\int_{\mathbb R}g\ne0\)，结论是否仍成立？
\end{exercise}
\begin{solution}
两个传播量 \(p(t,x)=(g+cf')(x+ct)\)、\(q(t,x)=(g-cf')(x-ct)\) 分别支撑在以 \(-ct\)、\(ct\) 为中心、半径 \(R\) 的区间内；在给定时间它们不相交。因此
\[
 \|u_t\|_2^2-c^2\|u_x\|_2^2=\int p(t,x)q(t,x)\,\mathrm{d}x=0.
\]
再结合总能量定义即可。证明只用两个导数量的支集，未要求 \(\int g=0\)。初速度总积分非零时，中间仍可留下常位移平台，但平台的速度和空间梯度都为零，不贡献能量。
\end{solution}

% 自拟：时间冲量的分布跳跃，增加Duhamel积分边界的检验。
\Needspace{6\baselineskip}
\begin{exercise}\label{b2:ex:09-impulse}
设 \(c>0\)、\(\tau>0\)、\(\mathbb K\in\{\mathbb R,\mathbb C\}\)，\(h\in C_c^\infty(\mathbb R;\mathbb K)\)。对 \(s\geq0\)，令 \(S_c(s)h(x)=(2c)^{-1}\int_{x-cs}^{x+cs}h(y)\,\mathrm{d}y\)。定义 \(u(t)=0\) 于 \(t<\tau\)，\(u(t)=S_c(t-\tau)h\) 于 \(t\geq\tau\)。证明在 \(\mathcal D'(\mathbb R_t\times\mathbb R_x)\) 中
\[
 u_{tt}-c^2u_{xx}=\delta_\tau(t)\otimes h(x).
\]
指出位移和速度在 \(\tau\) 处的跳跃，并解释当 \(h\ne0\) 时为何这种解不属于任何跨越 \(\tau\) 的时间区间上的 \(C^1_tL^2_x\) 类。
\end{exercise}
\begin{solution}
两侧分别满足齐次方程，且 \(u(\tau+)=u(\tau-)=0\)，\(u_t(\tau+)=h\)、\(u_t(\tau-)=0\)。对紧支集测试函数在 \(t=\tau\) 两侧分别积分分部，位移连续使一次时间分布导数不含点质量；再次求导产生速度跳跃 \(h\) 乘 \(\delta_\tau\)。空间导数不跨时间界面，没有额外项。故所求等式成立。除非 \(h=0\)，速度在 \(L^2\) 中不连续；外力也不是 \(L^1_tL^2_x\) 函数，所以不存在与能量解定理矛盾的问题。
\end{solution}

% 自拟：散度型变系数的能量与最大速度，不套用常速显式公式。
\Needspace{6\baselineskip}
\begin{exercise}\label{b2:ex:09-variable-speed}
设 \(a_0,A>0\)，\(a\in C^1(\mathbb R;\mathbb R)\) 满足 \(a_0\leq a(x)\leq A\)，且 \(\mathbb K\in\{\mathbb R,\mathbb C\}\)。设 \(u\in C^2([0,\infty)\times\mathbb R;\mathbb K)\) 满足 \(u_{tt}-\partial_x(a(x)u_x)=0\)，并假设对每个 \(T>0\)，存在 \(R_T>0\) 使 \(u(t,x)=0\) 于 \(0\leq t\leq T\)、\(|x|>R_T\)。证明能量
\[
 E_a(t)=\frac12\int_{\mathbb R}(|u_t|^2+a(x)|u_x|^2)\,\mathrm{d}x
\]
守恒，并用收缩区间证明传播速度不超过 \(C=\sqrt A\)：若 \(t_0>0\)、\(x_0\in\mathbb R\)，初始位移和初始速度在 \([x_0-Ct_0,x_0+Ct_0]\) 上均为零，则 \(u\) 在 \(0\leq t\leq t_0\)、\(|x-x_0|\leq C(t_0-t)\) 上为零。说明本题的支集条件只涉及有限时间带中的空间变量。
\end{exercise}
\begin{solution}
记 \(e=(|u_t|^2+a|u_x|^2)/2\)。乘以 \(\overline{u_t}\) 取实部得到 \(\partial_te=\partial_x\operatorname{Re}(a u_x\overline{u_t})\)；系数与时间无关，故不出现 \(a_t\) 项。在任意有限时间带上，共同的紧空间支集消去无穷远边界项，积分即得守恒。记 \(C=\sqrt A\)，在速度 \(C\) 收缩的区间上，任一端点的通量绝对值满足
\[
 |\operatorname{Re}(a u_x\overline{u_t})|
 \leq\sqrt a\,\dfrac{|u_t|^2+a|u_x|^2}{2}\leq C\cdot e.
\]
具体地，令 \(\alpha(t)=x_0-C(t_0-t)\)、\(\beta(t)=x_0+C(t_0-t)\)，则
\[
 \frac{\mathrm{d}}{\mathrm{d}t}\int_{\alpha(t)}^{\beta(t)}e(t,x)\,\mathrm{d}x
 =\operatorname{Re}(a u_x\overline{u_t})(t,\beta)-C e(t,\beta)
 -\operatorname{Re}(a u_x\overline{u_t})(t,\alpha)-C e(t,\alpha)\leq0.
\]
零初始数据也使底边上的空间导数为零，故局部能量从零开始并始终为零。由 \(a\geq a_0>0\)，内部的速度和梯度均为零；沿三角形内的竖直线回到初值线，得到 \(u(t,x)=u(0,x)=0\)，边界与顶点由连续性补齐。这一局部结论只使用有限三角形，不需要无穷远支集条件；共同紧空间支集用于前面的总能量守恒，不能读作整个时空的紧支集。此法给出统一速度上界，并未断言每条实际传播曲线都以 \(\sqrt A\) 运动。
\end{solution}

% 自拟：低频扩张族说明齐次能量不能代替非齐次H1范数。
\Needspace{6\baselineskip}
\begin{exercise}\label{b2:ex:09-missing-zero-order}
设 \(c>0\)、\(\mathbb K\in\{\mathbb R,\mathbb C\}\)，取非零 \(\eta\in C_c^\infty(\mathbb R;\mathbb K)\)，对正整数 \(n\) 令 \(f_n(x)=n^{-1/2}\eta(x/n)\)、\(g_n=0\)。令 \(u_n\) 为 \(\mathbb R_x\) 上满足 \(u_{n,tt}-c^2u_{n,xx}=0\)、\(u_n(0)=f_n\)、\(u_{n,t}(0)=g_n\) 的经典解，记 \(E_n(t)=\frac12(\|u_{n,t}(t)\|_2^2+c^2\|u_{n,x}(t)\|_2^2)\)。证明初始能量趋零，而 \(f_n\) 不趋于零于 \(L^2\)。对每个 \(T>0\)，证明 \(\sup_{0\leq t\leq T}\|u_n(t)-f_n\|_2\to0\)，并解释能量稳定性为何必须单独保留初始位移的零阶误差。
\end{exercise}
\begin{solution}
换元给出 \(\|f_n\|_2=\|\eta\|_2\)、\(\|f_n'\|_2=n^{-1}\|\eta'\|_2\)，所以 \(E_n(0)=c^2\|\eta'\|_2^2/(2n^2)\to0\)。能量守恒控制 \(\|u_{n,t}(t)\|_2\leq c\|\eta'\|_2/n\)，积分后得到
\[
 \sup_{0\leq t\leq T}\|u_n(t)-f_n\|_2\leq\dfrac{cT}{n}\|\eta'\|_2\longrightarrow0.
\]
因此 \(\|u_n(t)\|_2\) 在有限时间上一致趋于非零常数 \(\|\eta\|_2\)。几乎不弯曲而分布在越来越大区域上的位移，能量很小却仍有固定平方质量；只比较梯度和速度不能恢复 \(H^1\) 范数控制。
\end{solution}

